\begin{bt}%[1H4H1-3]
% nguon: Archive/SP_DA_SauLuoi/1-TL-THPT-NamKyKhoiNghia-GHKI-NH23-24.tex #2449
%[Dự án đề kiểm tra Toán 11 GHKI NH23-24 - Quan Ón]%[THPT Nam Kỳ Khởi Nghĩa - Tp HCM]
	Cho hình chóp $S.ABCD$ có đáy $ABCD$ là hình bình hành tâm $O$; $G$ là trọng tâm tam giác $SCD$; $H$ là trọng tâm tam giác $ACD$.
	\begin{enumerate}
		\item Tìm giao tuyến của hai mặt phẳng $(SAC)$ và $(SBD)$.
		\item Chứng minh hai đường thẳng $HG$ và $SA$ song song.
		\item Tìm giao điểm của đường thẳng $OG$ với mặt phẳng $(SAD)$.
	\end{enumerate}
	\loigiai{
		\begin{center}
			\begin{tikzpicture}[>=stealth,line join=round,line cap=round,font=\footnotesize,scale=0.9]
				\path 
				(0.57,3.34) coordinate (S)
				(0,0) coordinate (A)
				(-2,-2) coordinate (B)
				(3,-2) coordinate (C)
				(5,0) coordinate (D)
				(intersection of A--C and B--D) coordinate (O)
				($(C)!0.5!(D)$) coordinate (M)
				($(S)!0.5!(C)$) coordinate (N)
				($(A)!2/3!(M)$) coordinate (H)
				($(S)!2/3!(M)$) coordinate (G)
				($(D)-(O)+(N)$) coordinate (a)
				($(a)!1.6!(D)$) coordinate (x)
				(intersection of D--a and O--G) coordinate (E);
				\draw (S)--(B)--(C)--(D)--(S)--(C) (S)--(M) (D)--(N) (G)--(E)--(x);
				\draw[dashed] (B)--(A)--(D)--(B) (C)--(A)--(S)--(O)--(N) (A)--(M) (H)--(G)--(O);
				\foreach \l/\g in {A/180,B/-135,C/-45,D/0,O/-90,S/90,N/180,M/-45,G/85,H/-70,E/0}
				\draw[fill=black] (\l) circle (1pt) +(\g:.3) node{$\l$};
				\fill (x) node[right]{$x$};
			\end{tikzpicture}
		\end{center}
		\begin{enumerate}
			\item Ta có $S \in (SAC)\cap (SBD)$. $\hfill (1)$\\
			Vì $O$ là tâm của hình bình hành $ABCD$ nên $O \in AC$, $O \in BD$.\\
			Suy ra $\heva{&O \in AC,AC \subset (SAC)\\&O \in BD, BD \subset (SBD)} \Rightarrow O \in (SAC)\cap (SBD)$. $\hfill (2)$\\
			Từ $(1)$ và $(2)$, suy ra $SO = (SAC)\cap (SBD)$.
			\item Gọi $M$, $N$ lần lượt là trung điểm của các cạnh $CD$ và $SC$.\\
			Vì $G$ là trọng tâm tam giác $SCD$ nên $\dfrac{MG}{MS} = \dfrac{1}{3}$.\\
			Mặt khác, ta có $H$ là trọng tâm tam giác $ACD$ suy ra $\dfrac{MH}{MA} = \dfrac{1}{3}$.\\
			Do đó $\dfrac{MG}{MS} = \dfrac{MH}{MA} = \dfrac{1}{3} \Rightarrow HG \parallel SA$.
			\item Vì $G$ là trọng tâm tam giác $SCD$ nên $\dfrac{DG}{DN} = \dfrac{2}{3}$.\\
			Vì $H$ là trọng tâm tam giác $ACD$ nên $\dfrac{DH}{DO} = \dfrac{2}{3}$.\\
			Do đó $\dfrac{DG}{DN} = \dfrac{DH}{DO} = \dfrac{2}{3} \Rightarrow ON \parallel HG$.\\
			Mà $HG \parallel SA$ (cmt) nên $ON \parallel SA$.\\
			Xét $OG \subset (OGD)$.\\
			Ta có $\heva{&D \in (SAD)\cap (OGD)\\
			&ON \subset (OGD)\\&SA \subset (SAD)\\&ON \parallel SA \textrm{ (cmt)}} \Rightarrow (OGD)\cap (SAD) = Dx\parallel ON\parallel SA$.\\
			Trong $(OGD)$, gọi $E = OG \cap Dx$.\\
			Khi đó $\heva{&E \in OG\\&E \in Dx, Dx \subset (SAD)} \Rightarrow E = OG \cap (SAD)$.
		\end{enumerate}
	}
\end{bt}

\begin{bt}%[1H4V3-3]
% nguon: Archive/SP_DA_SauLuoi/1-TL-TN-THPT-NGUYENCHITHANH-GHKI-NH23-24.tex #2502
%[Dự án đề kiểm tra Toán 11 GHKI NH23-24- Nguyễn Sơn]%[THPT Nguyễn Chí Thanh - Tp HCM]
Cho hình chóp $ S.ABCD $, $ ABCD $ là hình thang có đáy là $ AD $ và $ BC $, $ AD = 2BC $.
Gọi $ E $ là trung điểm $ SA $, $ M $ là trọng tâm $ \triangle SAD $, $ G $ là giao điểm của $ AC $ và $ BD $.
\begin{listEX}[1]
	\item Tìm giao tuyến của hai mặt phẳng $ (MBC) $ và $ (SAD) $.
	\item Tìm $ F $ là giao điểm của $ MC $ và mặt phẳng $ (SBD) $.
	\item Chứng minh $ MG $ song song với $ BE $.
\end{listEX}
\loigiai{
\begin{center}
\begin{tikzpicture}[
	>=stealth, line join=round, line cap=round, 
	scale=1,
	declare function={r=3; a=75;}]
% ve diem
	\path 
		(0,0) coordinate (A)
		(-3,-3) coordinate (B)
		(2,-3) coordinate (C)
		(10,0) coordinate (D)
		(-1,6) coordinate (S)
		($ (S)!.5!(A) $) coordinate (E)
		($ (E)!{1/3}!(D) $) coordinate (M)
		(intersection of A--C and B--D) coordinate (G)
		($ (S)!{2/3}!(A) $) coordinate (H)
		($ (S)!{2/3}!(D) $) coordinate (K)
		(intersection of B--K and M--C) coordinate (F)
		;
	\draw[thick] 
		(S)--(B)--(C)--(D)--(S)--(C)
		;
	\draw[dashed]
		(S)--(A)--(B)--(D)--(E)
		(D)--(A)--(C)--(M)--(B)
		(H)--(K)--(B);
% phan giac goc A
%	\path[name path=t1] (A) -- (B) -- (C) -- cycle;
%	\path[name path=c1] (A) circle (2);
%	\path[name intersections={of = c1 and t1, by={a,b}}];
%	\path 
%		($ (a)!.5!(b) $) coordinate (d)
%		(intersection of A--d and B--C) coordinate (D);
% ve duong tron
%	\draw [name path=dtO]  (O) circle (r);
% giao diem cua duong tron voi duong tron/duong thang
%	\path[name intersections={of= dtO and dtI}] coordinate (B) at (intersection-1)  coordinate (C) at (intersection-2); 
% hinh chieu vuong goc cua B len AC
%	\path ($(A)!(B)!(C)$) coordinate (E);
% giao diem cua doan thang voi doan thang
%	\path (intersection of O--S and C--B) coordinate (M);
% danh dau goc
	\foreach \a/\b/\c in {}	{
		\draw pic [draw=black,angle radius = 10] {right angle = \a--\b--\c};
	}
%	\draw pic ["$1$", draw, angle eccentricity=1.3, angle radius=20] {angle = B--A--C};
% danh dau diem
	\foreach \x/\g in {S/90, A/135, B/-135, C/-45, D/0, E/135, M/45, G/-90, H/180, K/45, F/135}{
		\draw[fill=black] (\x) circle (1.5pt)+(\g:.4) node{$\x$};}
\end{tikzpicture}
\end{center}
\begin{listEX}[1]
	\item Xét hai mặt phẳng $ (MBC) $ và $ (SAD) $, ta thấy
	$
		\heva{
			&M \in (MBC)\cap (SAD)\\
			&BC \subset (MBC)\\
			&AD \subset (SAD)\\
			&AD \parallel BC
		}
	$
	\\
	$ \Rightarrow $ giao tuyến của $ (MBC) $ và $ (SAD) $ là đường thẳng qua $ M $, song song với $ AD $ và $ BC $, cắt $ SA $ và $ SD $ lần lượt tại $ H $ và $ K $.
	\item Trong mặt phẳng $ (MBC) $, gọi $ F $ là giao điểm của $ BK $ và $ MC $.
	Ta thấy	$\heva{&F \in MC \\ &F \in BK \subset (SBD) }	$
	\\
	Suy ra $ F $ là giao điểm của $ MC $ và mặt phẳng $ (SBD) $.
	\item $ \triangle GBC $ đồng dạng $ \triangle GDA $ cho ta
	$\dfrac{GB}{GD} = \dfrac{BC}{AD} = \dfrac{1}{2} \quad (1)$.\\
	$ M $ là trọng tâm $ \triangle SAD $ nên $\dfrac{ME}{MD} = \dfrac{1}{2} \quad (2)$.\\
	Từ $ (1) $ và $ (2) $ suy ra $ \dfrac{GB}{GD} = \dfrac{ME}{MD} $. 
	\\
	Do vậy theo định lý Thalet đảo, ta có $ MG \parallel BE$.
\end{listEX}
}
\end{bt}

\begin{bt}%[1H4V1-3]
% nguon: Archive/SP_DA_SauLuoi/1-TL-TN-THPT-PHAMPHUTHU-GHKI-NH23-24.tex #2597
%[Dự án đề kiểm tra Toán 11 GHKI NH23-24- Nguyễn Sơn]%[THPT Phạm Phú Thứ - Tp HCM]
Cho hình chóp $ S.ABCD $ có đáy $ ABCD $ là hình bình hành tâm $O$. Gọi $M, N$ lần lượt là trung điểm của $SB, SD$. Trên đoạn thẳng $SC$ lấy điểm $P$ sao cho $SP > PC$. Biết $MN$ cắt $SO$.
\begin{listEX}[1]
	\item Tìm giao tuyến của hai mặt phẳng $ (SAD) $ và $ (SBC) $.
	\item Tìm giao tuyến của hai mặt phẳng $ (SAC) $ và $(MNP)$.
	\item Tìm giao điểm $Q$ của đường thẳng $SA$ và mặt phẳng $(MNP)$.
	\item Gọi $I,J$ lần lượt là giao điểm của $QM$ và $AB, QP$ và $AC$. Chứng minh $IJ, QN$ và $AD$ đồng quy.
\end{listEX}
\loigiai{
\begin{center}
\begin{tikzpicture}[
	>=stealth, line join=round, line cap=round, 
	scale=0.7,
	declare function={r=3; a=75;}]
% ve diem
	\path 
		(0,0) coordinate (A)
		(-3,-3) coordinate (D)
		(6,-3) coordinate (C)
		(9,0) coordinate (B)
		(2,6) coordinate (S)
		(-3,1) coordinate (x1)
		(4,8) coordinate (x)
		($ (S)!.5!(B) $) coordinate (M)
		($ (S)!.5!(D) $) coordinate (N)
		($ (S)!.85!(C) $) coordinate (P)
		(intersection of A--C and B--D) coordinate (O)	
		(intersection of S--O and M--N) coordinate (K)
		(intersection of P--K and S--A) coordinate (Q)
		(intersection of Q--M and A--B) coordinate (I)	
		(intersection of Q--P and A--C) coordinate (J)	
		(intersection of I--J and A--D) coordinate (E)	
		;
	\draw[thick] (S)--(B)--(C)--(D)--(S)--(C) (x1)--(x) (M)--(P) (N)--(P) (B)--(I)--(M) (P)--(J)--(C);
	\draw[dashed](S)--(A)--(B)--(D)	(D)--(A)--(C) (S)--(O) (M)--(N) (P)--(Q) (P)--(J)--(C) (P)--(Q)--(M) (Q)--(N);
	\draw  (I)--(E)--(A) (E)--(Q);
	\foreach \a/\b/\c in {}	{
		\draw pic [draw=black,angle radius = 10] {right angle = \a--\b--\c};	}
% danh dau diem
	\foreach \x/\g in {S/90, A/135, B/-90, C/-140, D/180, M/90, N/180, O/-90, P/0, x/90, K/-120, Q/20, I/-90, J/-90, E/-90}{
		\draw[fill=black] (\x) circle (1.5pt)+(\g:.4) node{$\x$};}
\end{tikzpicture}
\end{center}
\begin{listEX}[1]
	\item Ta có $\heva{&S\in (SAD)\cap (SBC) \\ & AD\subset (SAD)\\&BC\subset(SBC)\\&AD\parallel BC}\Rightarrow (SAD)\cap (SBC) = Sx $ với $Sx \parallel AD \parallel BC$.
	\item Trong mặt phẳng $(SBD)$, gọi $K$ là giao điểm của $MN$ và $SO$.\\
	 Suy ra  $K\in (SAC)\cap (MNP)$.\\
	Ta lại có  $P\in (SAC)\cap (MNP)$.\\
	Suy ra $KP=(SAC)\cap(MNP)$.
	\item Trong mặt phẳng $(SAC)$ gọi $Q$ là giao điểm của $SA$ và $KP$.\\
	Suy ra $\heva{&Q\in SA\\&Q\in KP\subset (MNP)}\Rightarrow Q=SA\cap (MNP)$.
	\item Trong mặt phẳng $(SAD)$ gọi $E$ là giao điểm của $QN$ và $AD$.\\
	Ta cần chứng minh ba điểm $E, J, I $ thẳng hàng.\\
	Ta có $\heva{&E\in AD\subset (ABCD)\\&E\in QN\subset(MNP)}\Rightarrow E\in (ABCD)\cap (MNP)$\quad (1).\\
	Ta có $\heva{&J\in AC\subset (ABCD)\\&J\in QP\subset(MNP)}\Rightarrow J\in (ABCD)\cap (MNP)$\quad (2).\\
	Ta có $\heva{&I\in AB\subset (ABCD)\\&I\in QM\subset(MNP)}\Rightarrow I\in (ABCD)\cap (MNP)$\quad (3).\\
	Từ (1) (2) và (3) suy ra $E, I, J$ cùng thuộc giao tuyến hai mặt phẳng $(ABCD)$ và $(MNP)$. Suy ra $E, J, I $ thẳng hàng. Vậy ba đường thẳng $IJ, QN$ và $AD$ đồng quy.
\end{listEX}
}
\end{bt}

\begin{bt}%[1H4H1-3]
% nguon: Archive/SP_DA_SauLuoi/1-TL-THPT-TranKhaiNguyen-GHKI-NH23-24.tex #2631

	Cho hình chóp $S.ABCD$ có đáy là hình bình hành tâm $O$. Gọi $M$, $N$ lần lượt là trung điểm của $SB$, $SD$. Lấy $P$ thuộc đoạn $SC$ và không là trung điểm của $SC$.
	\begin{enumerate}
		\item Tìm giao tuyến của hai mặt phẳng $(SAB)$ và $(SBD)$. Tìm giao điểm $E$ của đường thẳng SO và mặt phẳng $(MNP)$.
		\item Chứng minh $MN\parallel BD$. Từ đó tìm giao tuyến của hai mặt phẳng $(ABCD)$ và $(CMN)$.
	\end{enumerate}
	\loigiai
	{
		\begin{center}
			\begin{tikzpicture}[>=stealth,line join=round,line cap=round,font=\footnotesize,scale=1]
				\coordinate (B) at (0,0);
				\coordinate (A) at (1,1.2);
				\coordinate (C) at (3.5,0);
				\coordinate (D) at ($(C)-(B)+(A)$);
				\coordinate (S) at ($(A)+(100:3.5)$);
				\coordinate (M) at ($(S)!0.5!(B)$);
				\coordinate (N) at ($(S)!0.5!(D)$);
				\coordinate (P) at ($(S)!0.6!(C)$);
				\coordinate (O) at ($(A)!0.5!(C)$);
				\draw ($(C)+(N)-(M)$) node[below]{$d$}--(C)--($(C)+(M)-(N)$);
				\draw (B)--(C)--(D)--(S)--cycle (S)--(C) (M)--(P)--(N)--(C)--cycle;
				\draw[dashed] (A)--(D) (S)--(A)--(B) (A)--(C) (B)--(D) (M)--(N) (S)--(O);
				\foreach \diem/\goc in {A/180,B/-90,C/-90,D/0,S/90,M/180,P/0,N/0,O/-90} \fill[black](\diem) circle (1pt) ($(\diem)+(\goc:3mm)$) node{$\diem$};
			\end{tikzpicture}
		\end{center}
		\begin{enumerate}
			\item Xét hai mặt phẳng $(SAB)$ và $(SBD)$. Ta có
			\begin{align*}
				&\heva{&S\in(S A B)\\ &S\in(S B D)}\Rightarrow S\in(S A B)\cap(S B D)\\ 
				&\heva{&B\in(S A B)\\ &B\in(S B D)}\Rightarrow B\in(S A B)\cap(S B D)
			\end{align*}
			Vậy $(SAB)\cap (SBD)=SB$.\\
			Trong mặt phẳng $(SBD)$, gọi $E=SO\cap MN$.\\
			Ta có $\heva{& E\in SO \\ & E\in MN,MN\subset (MNP)}\Rightarrow \heva{& E\in SO \\ & E\in (MNP)}\Rightarrow E=SO\cap (MNP)$.\\
			Vậy $E$ là giao điểm của $SO$ và $(MNP)$.
			\item Xét tam giác $SBD$ có $M$ là trung điểm $SB$, $N$ trung điểm $SD$, suy ra $MN$ là đường trung bình của tam giác $SBD\Rightarrow MN\parallel BD$.\\
			Ta có $\heva{& C\in (ABCD)\cap (CMN) \\ & MN\parallel BD \\ & MN\subset (CMN), BD\subset (ABCD)}\Rightarrow (ABCD)\cap (CMN)=d$ với $d\parallel MN\parallel BD$ và $d$ đi qua $C$.
		\end{enumerate}
	}
\end{bt}

\begin{bt}%[1H4C5-3]
% nguon: Archive/SP_KhoLatex/ID6_old/4.0.1.tex #2857

		Cho hình hộp chữ nhật $ABCD\cdot A'B'C'D'$. Xác định các điểm $M,N$ tương ứng trên các đoạn $AC'$, $B'D'$ sao cho $MN$ song song với $BA'$ và tính tỉ số $\dfrac{MA}{MC'}$.
\loigiai{
	\begin{center}
	\begin{tikzpicture}[>=stealth,line join=round,line cap=round,font=\footnotesize,scale=1]
		\coordinate (A') at (0,0);
		\coordinate (D') at (5,0);
		\coordinate (B') at (1,-2);
		\coordinate (C') at ($(B')+(D')-(A')$);
		\coordinate (A) at ($(A')+(2,3)$);
		\coordinate (B) at ($(B')+(2,3)$);
		\coordinate (C) at ($(C')+(2,3)$);
		\coordinate (D) at ($(D')+(2,3)$);
		\coordinate (K) at ($(B')!2!(A')$);
		\coordinate (N) at (intersection of C'--K and B'--D');
		\coordinate (M) at ($(C')!1/3!(A)$);
		\draw (A')--(B')--(C') (A')--(A)--(D)--(C)--(B)--(A) (B)--(B') (C)--(C') (A')--(K)--(A) ;
		\draw[dashed] (C')--(D')--(A') (D)--(D')--(B') (A)--(C')--(K) (N)--(M);
		\foreach \i/\g in {A'/180,B'/180,C'/0,D'/0,A/90,B/180,C/0,D/0,K/180,N/-90,M/-90}{\draw[fill=black](\i) circle (1.5pt) ($(\i)+(\g:3mm)$) node[scale=1]{$\i$};}
	\end{tikzpicture}
\end{center}
		Xét phép chiếu song song lên mặt phẳng $\left(A'B'C'D'\right)$ theo phương $BA'$, ta có \\
		$N$ là hình chiếu của $M$ hay $N$ là giao điểm của $B'D'$ với hình chiếu của $AC'$ qua phép chiếu này.\\
		Do đó ta xác định được $M,N$ như sau:
		\begin{itemize}
			\item Trên tia $B'A'$, lấy điểm $K$ sao cho $A'K=B'A'$ thì $ABA'K$ là hình bình hành nên $AK\parallel BA'$, suy ra $K$ là hình chiếu của $A$ qua phép chiếu song song.
			\item Gọi $N=B'D'\parallel KC'$. Đường thẳng qua $N$ và song song với $AK$ cắt $AC'$ tại $M$.
		\end{itemize}
		Từ đó ta xác định được điểm $M,N$ cần tìm.\\
		Theo định lý Ta-let, ta có $\dfrac{MA}{MC'}=\dfrac{NK}{NC'}=\dfrac{KB'}{C'D'}=2$.}
\end{bt}

\begin{bt}%[1H4C5-3]
% nguon: Archive/SP_KhoLatex/ID6_old/4.0.1.tex #2858

		Cho hình hộp chữ nhật $ABCD\cdot A'B'C'D'$. Gọi $M,N$ lần lượt là trung điểm $CD$ và $DD'$. Xác định đường thẳng $d$ đi qua $M$ đồng thời cắt $BN$ và $AB'$.
	\loigiai{
\begin{center}
	\begin{tikzpicture}[>=stealth,line join=round,line cap=round,font=\footnotesize,scale=1]
		\coordinate (A) at (0,0);
		\coordinate (D) at (6,0);
		\coordinate (B) at (-2,-2);
		\coordinate (C) at ($(B)+(D)-(A)$);
		\coordinate (A') at ($(A)+(-1,5)$);	
		\coordinate (B') at ($(B)+(-1,5)$);	
		\coordinate (C') at ($(C)+(-1,5)$);	
		\coordinate (D') at ($(D)+(-1,5)$);	
		\coordinate (M) at ($(C)!.5!(D)$);
		\coordinate (N) at ($(D')!.5!(D)$);
		\coordinate (J) at ($(A)!.5!(B')$);
		\coordinate (I') at ($(A)!.5!(M)$);
		\coordinate (I) at ($(J)!.5!(M)$);
		\coordinate (N') at (intersection of B--I' and C--D);
		\draw(B')--(C')--(D')--(A')--(B')--(B)--(C)--(D)--(D') (C)--(C') (N)--(N')--(D);
		\draw[dashed] (A')--(A)--(D) (A)--(B)--(N) (B')--(A)--(M)--(J) (I)--(I') (B)--(N') ;
		\foreach \i/\g in {A/180,B/180,C/-90,D/0,A'/90,B'/180,C'/0,D'/0,M/0,N/0,J/80,I'/-90,I/90,N'/0}{\draw[fill=black](\i) circle (1.5pt) ($(\i)+(\g:3mm)$) node[scale=1]{$\i$};}
	\end{tikzpicture}
\end{center}
		Giả sử ta đã dựng được đường thẳng $d$ cắt $BN$ và $AB'$. Gọi $I$, $J$ lần lượt là giao điểm của $d$ với $BN$ và $AB'$.\\
		Xét phép chiếu song song theo phương chiếu $B'A$ lên mặt phẳng $\left(ABCD\right)$. Khi đó, ta có hình chiếu của các điểm $J$, $I$, $M$ lần lượt là $A$, $I'$, $M$.\\
		Do $J$, $I$, $M$ thẳng hàng nên $A$, $I'$, $M$ cũng thẳng hàng.\\
		Gọi $N'$ là hình chiếu của $N$ thì $BN'$ là hình chiếu của $BN$.\\
		Vì $I\in BN\Rightarrow I'\in BN'\Rightarrow I'=AM\cap BN'$\\
		Từ những phân tích trên, ta suy ra cách dựng đường thẳng $d$ như sau :
		\begin{itemize}
			\item Lấy $I'=BN'\cap AM$.
			\item Trong $\left(BNN'\right)$ dựng $II'\parallel NN'$ cắt $BN$ tại $I$.
			\item Vẽ đường thẳng $MI$, đó cũng là đường thẳng $d$ cần tìm.
		\end{itemize}}
\end{bt}

\begin{bt}%[1H4H1-4]
% nguon: Archive/SP_KhoLatex/ID6_old/1-CK1-OT-De-25.tex #2878

	Cho hình chóp $S. A B C D$ có đáy $A B C D$ là hình bình hành tâm $O$. Điểm $M$ thuộc cạnh $B C$ sao cho $B M=2 M C$. Tìm giao điểm của đường thẳng $A B$ và mặt phẳng $(S D M)$,
	\loigiai{
		\immini{
			Trong $(A B C D)$, $A B$ cắt $D M$ tại $N$.\\
			$\heva{&N \in A B \\&N \in D M,\, D M \subset(S D M).}$\\
			Vậy $ N$ là giao điểm của $A B$ và $(S D M)$.
		}{
			\begin{tikzpicture}[line join = round, line cap = round,>=stealth,font=\footnotesize,scale=.7]
				\coordinate[label=left:$B$] (A) at (0,0);
				\coordinate[label=right:$C$] (B) at ($(A)+(6,0)$);
				\coordinate[label=left:$A$] (D) at ($(A)+(1.5,2)$);
				\coordinate[label=right:$D$] (C) at ($(B)+(D)-(A)$);
				\coordinate[label=above:$S$] (S) at (.5,6);
				\coordinate[label=below:$M$] (M) at ($(A)!2/3!(B)$); % Định nghĩa điểm M thỏa mãn \vt{AM}=1/2*\vt{AB}
				\coordinate[label=below:$N$] (N) at ($(D)!3!(A)$);
				\draw (S)--(C)--(B)--(A)--(S)--cycle (S)--(B) (A)--(N) (N)--(M);
				\draw[dashed] (S)--(D)--(C) (A)--(D)  (C)--(M);
				\fill[black] (A) circle (1.5pt);
				\fill[black] (B) circle (1.5pt);
				\fill[black] (C) circle (1.5pt);
				\fill[black] (D) circle (1.5pt);
				\fill[black] (S) circle (1.5pt);
				\fill[black] (M) circle (1.5pt);
				\fill[black] (N) circle (1.5pt);
			\end{tikzpicture}
	}}
\end{bt}

\begin{bt}%[1H4V3-2]
% nguon: Archive/SP_KhoLatex/ID6_old/1-CK1-OT-De-25.tex #2880

	Cho hai hình vuông $ABCD$, $ABEF$ nằm trong hai mặt phẳng khác nhau. Trên các đường chéo $AC$ và $BF$ ta lấy các điểm $M$, $N$ sao cho $AM=BN$. Chứng minh rằng $MN\parallel \left(DEF\right)$.
	\loigiai{
		\immini{Trong $\left(ABCD\right)$ qua $M$ kẻ $MM'\parallel AB$ $\left(M'\in AD\right)$.\\
			Trong $\left(ABEF\right)$ qua $N$ kẻ $NN'\parallel AB$ $\left(N'\in AF\right)$.\\
			Ta có
			$\heva{
				&\dfrac{AM'}{AD}=\dfrac{AM}{AC}\\
				&\dfrac{AN'}{AF}=\dfrac{BN}{BF}\\
				&AM=BN;\, AC=BF}$\\
			$\Rightarrow\dfrac{AM'}{AD}=\dfrac{AN'}{AF}\Rightarrow M'N'\parallel DF$.\\
			Ta có $NN'\parallel AB\parallel EF\Rightarrow\left(MM'N'N\right)\parallel \left(DEF\right)$.\\
			Mà $MN\subset\left(MM'N'N\right) \Rightarrow MN\parallel \left(DEF\right)$.}{
			\begin{tikzpicture}[line join = round, line cap = round,>=stealth,font=\footnotesize,scale=.7]
				\coordinate[label=left:$A$] (A) at (0,0);
				\coordinate[label=right:$D$] (D) at ($(A)+(7,0)$);
				\coordinate[label=left:$B$] (B) at ($(A)+(2,2)$);
				\coordinate[label=right:$C$] (C) at ($(D)+(B)-(A)$);
				\coordinate[label=above:$E$] (E) at ($(B)+(0,7)$);
				\coordinate[label=above:$F$] (F) at ($(A)+(0,7)$);
				\coordinate[label=above:$M$] (M) at ($(A)!1/3!(C)$);
				\coordinate[label=above right:$N$] (N) at ($(B)!1/3!(F)$);
				\coordinate[label=right:$N'$] (N1) at ($(A)!1/3!(F)$);
				\coordinate[label=below:$M'$] (M1) at ($(A)!1/3!(D)$);
				\draw (D)--(E)--(F)--cycle (E)--(C)--(D)--(A)--(F) (M1)--(N1);
				\draw[dashed] (A)--(B)--(C) (B)--(E) (A)--(C) (B)--(F) (M)--(N) (M)--(M1) (N1)--(N);
				\fill[black] (A) circle (1.5pt);
				\fill[black] (B) circle (1.5pt);
				\fill[black] (C) circle (1.5pt);
				\fill[black] (D) circle (1.5pt);
				\fill[black] (E) circle (1.5pt);
				\fill[black] (F) circle (1.5pt);
				\fill[black] (M) circle (1.5pt);
				\fill[black] (N) circle (1.5pt);
				\fill[black] (M1) circle (1.5pt);
				\fill[black] (N1) circle (1.5pt);
			\end{tikzpicture}
	}}
\end{bt}

\begin{bt}%[1H4V3-2]
% nguon: Archive/SP_KhoLatex/ID6_old/1-CK1-OT-De-31.tex #2897

	Cho hình chóp $S.ABCD$ có đáy $ABCD$ là một hình bình hành. Gọi $G$ là trọng tâm tam giác $SAB$, $I$ là trung điểm của $AB$ và $M$ là điểm trên cạnh $AD$ sao cho $AM=\dfrac{1}{3}AD$.
	\begin{enumerate}
		\item Đường thẳng đi qua $M$ và song song với $AB$ cắt $CI$ tại $N$. Chứng minh $NG\parallel \left( SCD \right)$.
		\item Chứng minh $MG\parallel \left( SCD \right)$.
	\end{enumerate}
	\loigiai{
\immini{	\begin{enumerate}
	\item Ta có $\heva{& \dfrac{IN}{IC}=\dfrac{BJ}{BC}=\dfrac{AM}{AD}=\dfrac{1}{3}\\&\dfrac{IG}{IS}=\dfrac{1}{3}}$\\
	$\Rightarrow \dfrac{IN}{IC}=\dfrac{IG}{IS}\Rightarrow NG\parallel SC$.\\
	Mà $SC\subset \left( SCD \right)$$\Rightarrow NG\parallel \left( SCD \right)$.
	\item 
	Kẻ $MJ\parallel AB$ và cắt $BC$ tại $J$.\\
	Gọi $E$ là giao điểm của $IM$ và $CD$.\\
	Ta có 			
	$\dfrac{IM}{IE}=\dfrac{AM}{AD}=\dfrac{1}{3}\Rightarrow \dfrac{IM}{IE}=\dfrac{IG}{IS}$\\
	$\Rightarrow MG\parallel SE$.\\
	Mà $SE\subset \left( SCD \right)\Rightarrow MG\parallel \left( SCD \right)$.
\end{enumerate}	}
{\begin{tikzpicture}[scale=0.7]
	\def\a{6}
	\def\h{4}
	\path 	(0:0) coordinate (A)
			++(0:\a) coordinate (D)
			++(-130:\a/2) coordinate (C)
			($(A)+(C)-(D)$) coordinate (B)
			($(A)+(90:\h)$) coordinate (S)
			($(A)!0.5!(B)$) coordinate (I)
			($(S)!2/3!(I)$) coordinate (G)
			($(A)!1/3!(D)$) coordinate (M)
			($(B)!1/3!(C)$) coordinate (J)
			(intersection of I--C and M--J) coordinate (N)
			(intersection of I--M and D--C) coordinate (E)
			(intersection of E--I and D--S) coordinate (K)
			(intersection of A--C and B--D) coordinate (O);%giao điểm O
	\draw[dashed,thick] 	(B)--(A)--(D)	(A)--(S) (M)--(J) (C)--(I)--(S) (N)--(G) (M)--(E) (G)--(M)--(I);
	\draw[thick] 			(B)-- (C)--(D)
			(B)--(S)	(C)--(S)	(D)--(S) (K)--(E)--(D) (S)--(E);
	\foreach \x/\g in {A/60,B/-135,C/-45,D/-45,S/90,G/45,I/-90,M/-55,J/-90,N/-90,E/45}
			\fill[black] 	(\x) circle (1.5pt)
			($(\g:3mm)+(\x)$) node {$\x$};	
\end{tikzpicture}}}
\end{bt}

\begin{bt}%[1C1C5-4]
% nguon: Archive/SP_KhoLatex/ID6_old/1H1G5-4.tex #3002
%[HK1, Lê Quí Đôn - Đà Nẵng, 2018]%[Đức Nguyễn, 11EX1819]
	Dựng về phía ngoài tam giác nhọn $ABC$ các hình vuông $ABMN$, $ACPQ$. Gọi $O_1$, $O_2$ lần lượt là tâm của hình vuông $ABMN$, $ACPQ$ và $I$ là trung điểm của $BC$.
	\begin{enumerate}
		\item Chứng minh tam giác $IO_1O_2$ vuông cân tại $I$.
		\item Chứng minh $AI$ vuông góc với $NQ$ và $AI=\dfrac{1}{2}NQ$
	\end{enumerate}
	\loigiai{
			\begin{center}
			\begin{tikzpicture}[scale=1, font=\footnotesize,line join=round, line cap=round,>=stealth]
			\pgfmathsetmacro\a{4}
			\tkzDefPoint(0,0){B}
			\tkzDefShiftPoint[B](0:\a){C}
			\tkzDefShiftPoint[B](75:\a-0.5){A}
			\tkzDefPointBy[rotation = center B angle 90](A) \tkzGetPoint{M}
			\tkzDefPointBy[translation = from B to A](M) \tkzGetPoint{N}
			\tkzDefPointBy[rotation = center A angle 90](C) \tkzGetPoint{Q}
			\tkzDefPointBy[translation = from A to Q](C) \tkzGetPoint{P}
			\tkzDefMidPoint(B,C) \tkzGetPoint{I}
			\tkzInterLL(A,M)(B,N)    \tkzGetPoint{O1}
			\tkzInterLL(A,P)(C,Q)    \tkzGetPoint{O2}
			\tkzDefMidPoint(N,C) \tkzGetPoint{X}
			\tkzDefMidPoint(B,Q) \tkzGetPoint{Y}
			\tkzDefPointBy[rotation = center A angle 90](O2) \tkzGetPoint{Z}
			\tkzDrawPolygon(A,B,C)
			\tkzDrawSegments(B,M M,N N,A C,P P,Q Q,A N,B A,M C,Q A,P N,C Q,B O1,O2 O2,I O1,I N,Q)
			\tkzDrawSegments[dashed](A,I X,O2 I,Y Z,Y A,Z)
			\tkzLabelPoints[above](Z,Q,N)
			\tkzLabelPoints[below](M,B,I,C)
			\tkzLabelPoints[right](P)
			\tkzLabelPoint[above left](Y){\footnotesize $Y$}
			\tkzLabelPoint[left](O1){\footnotesize $O_1$}
			\tkzLabelPoint[above left](O2){\footnotesize $O_2$}
			\tkzLabelPoint[below](X){\footnotesize $X$}
			\tkzLabelPoints[above left](A)
		\tkzDrawPoints[fill=black](A,B,C,M,N,P,Q,X,Y,Z,I,O1,O2)
			\end{tikzpicture}
		\end{center}
		\begin{enumerate}
			\item Không mất tính tổng quát, giả sử $A$, $B$, $C$ có hướng như hình vẽ.\\
			Xét phép quay $\mathrm{Q}_{\left(A,\dfrac{\pi}{2}\right)}$ thì $N\mapsto B$, $C\mapsto Q$, suy ra $NC=BQ$ và $NC\perp BQ$.\\
			Từ giả thiết suy ra $IO_1$, $IO_2$ lần lượt là các đường trung bình trong các tam giác $BNC$, $CBQ$ nên suy ra $IO_1=\dfrac{1}{2}NC=\dfrac{1}{2}BQ=IO_2$ và $\widehat{O_1IO_2}=\dfrac{\pi}{2}$. Do đó, tam giác $IO_1O_2$ vuông cân tại $I$. 
			\item Gọi $X$, $Y$ lần lượt là  trung điểm của $NC$, $BQ$ và $Z=\mathrm{Q}_{\left(A,\dfrac{\pi}{2}\right)}\left(O_2\right)$.\\
			Qua $\mathrm{Q}_{\left(A,\dfrac{\pi}{2}\right)}$ thì $NC\mapsto BQ$ suy ra $X\mapsto Y$; $XO_2\mapsto YZ$.\\
			Do $IY$ là đường trung bình của tam giác $BCQ$ nên
			 $$IY\parallel QC, IY=\dfrac{1}{2}QC\Rightarrow
			  IY\parallel AZ, IY=AZ.$$
			Vậy $AIYZ$ là hình bình hành, suy ra $AI\parallel YZ$ và $AI=YZ\Rightarrow AI\perp XO_2$ và $AI=XO_2$.\\
			  Vì $XO_2$ là đường trung bình của tam giác $NCQ$ nên $AI\perp
			   NQ$ và $AI=\dfrac{1}{2}NQ$.
		\end{enumerate}
	}
\end{bt}

\begin{bt}%[1H4V1-4]
% nguon: Archive/SP_KhoLatex/ID6_old/GK1.2.tex #3273

	(1,0 điểm) Cho hình chóp tứ giác $S.ABCD$ có đáy $ABCD$ là hình thang $\left( BC//AD \right)$ và $G$ là trọng tâm của tam giác $SAD$.
	\begin{enumerate}
		\item Tìm giao điểm của đường thẳng $AC$ và mặt phẳng $\left( SBD \right)$.
		\item Tìm giao tuyến của mặt phẳng $\left( GBC \right)$ với các mặt bên của hình chóp đã cho. Đa giác tạo bởi các đoạn giao tuyến là hình gì?
	\end{enumerate}
	\loigiai{
		\begin{center}
			\begin{tikzpicture}[scale=0.9, font=\footnotesize,line join=round, line cap=round, >=stealth]
				\path
				(0:0) coordinate (A)
				(7,0) coordinate (D)
				(5,-2) coordinate (C)
				(2,-2) coordinate (B)
				(intersection of A--C and B--D)coordinate (O)+(85:5)coordinate (S)
				($(A)!0.5!(D)$) coordinate (M)
				($(S)!2/3!(M)$) coordinate (G)
				($(S)!2/3!(A)$) coordinate (E)
				($(S)!2/3!(D)$) coordinate (F)
				(intersection of S--O and A--M)coordinate (I)
				;
				\draw[dashed,thick]  (S)--(M) (A)--(C) (B)--(D) (A)--(M) (A)--(D) (F)--(E) (B)--(G)--(C);
				\draw[thick] (A)--(B)--(C)--(D) (S)--(A) (S)--(D) (S)--(C) (S)--(B) (E)--(B) (F)--(C) ;
				\foreach \x/\g in {A/185,B/-90,C/-90,D/-90,S/90,O/-90,M/-80,G/120,E/180,F/20}\draw[fill=white] (\x) circle (.03) +(\g:.33) node{$\x$};
			\end{tikzpicture}		
		\end{center}
		\begin{enumerate}
			\item Gọi $O$ là giao điểm của $AC$ và $BD$. Khi đó $O\in AC$ và $O\in BD$.\\
			Mà $BD\subset \left( SBD \right)\Rightarrow O\in \left( SBD \right)$ nên $O$ chính là giao điểm của đường thẳng $AC$ và mặt phẳng $\left( SBD \right)$.
			\item Mặt phẳng $\left( GBC \right)$ và mặt phẳng $\left( SAD \right)$ có điểm chung $G$, mặt khác $BC//AD$ nên giao tuyến của hai mặt phẳng là một đường thẳng đi qua $G$ và song song với $AD$.\\
			Đường thẳng này cắt $SA$ tại điểm $E$, cắt $SD$ tại điểm $F$.\\
			Vậy $\left( GBC \right)\cap \left( SBC \right)=BC$;\\
			$\left( GBC \right)\cap \left( SAB \right)=EB$;\\
			$\left( GBC \right)\cap \left( SAD \right)=EF$;\\
			$\left( GBC \right)\cap \left( SCD \right)=FC$;\\
			Do $EF//AD\Rightarrow EF//BC$ nên đa giác tạo bởi các đoạn giao tuyến nói trên là hình thang.
		\end{enumerate}
	}
\end{bt}

\begin{bt}%[1H4H2-2]
% nguon: Archive/SP_KhoLatex/ID6_old/1K4-12.tex #3285
%[Dự án 3 - TLDH - TeamTeXHoa - Võ Thị Thùy Trang]
    Cho hai hình bình hành $ABCD$ và $ABEF$ không cùng nằm trong một mặt phẳng. Gọi $O$ và $O'$ lần lượt là tâm của hai hình bình hành $ABCD$ và $ABEF$.
    \begin{enumerate}
        \item Chứng minh $OO'$ song song với các mặt phẳng $\left(ADF\right)$ và $\left(BCE\right)$.
        \item Gọi $G$ và $G'$ lần lượt là trọng tâm các tam giác $ABD$ và $ABF$. Chứng minh\break $GG' \parallel \left(DCEF\right)$.
    \end{enumerate}
    \loigiai{
        \immini{
            \begin{enumerate}
                \item Ta có $OO'$ là đường trung bình của tam giác $ACE$ và tam giác $BDF$ nên $OO' \parallel CE$ và $OO' \parallel DF$.\\
                Mà $CE\subset\left(BCE\right)$, $DF\subset\left(ADF\right)$ nên \break $OO'\parallel \left(BCE\right)$ và $OO'\parallel\left(ADF\right)$.
                \item Theo tính chất của trọng tâm tam giác, ta có 
        \end{enumerate}}
        {   \begin{tikzpicture}[scale=0.7,>=stealth, font=\footnotesize, line join=round, line cap=round]
                \coordinate[label=left:$A$] (A) at (2,1);   
                \coordinate[label=right:$B$] (B) at (7,1);  
                \coordinate[label=right:$C$] (C) at (9,-1); 
                \coordinate[label=left:$D$]  (D) at ($(A)+(C)-(B)$);
                \coordinate[label=right:$E$] (E) at (8,3);
                \coordinate[label=left:$F$]  (F) at ($(A)+(E)-(B)$);
                \coordinate[label=below:$O$] (O) at ($(A)!0.5!(C)$);
                \coordinate[label=above:$O'$] (O') at ($(A)!0.5!(E)$);
                \coordinate[label=below:$G$] (G) at ($(A)!2/3!(O)$);
                \coordinate[label=above:$G'$] (G') at ($(A)!2/3!(O')$);
                \draw (A)--(D)--(C) (A)--(F)--(E)--(C) (F)--(D);
                \draw[dashed] (A)--(B)--(C) (B)--(E) (A)--(C) (B)--(D) (A)--(E) (B)--(F) (O')--(O) (G)--(G'); 
                \foreach \diem in {A,B,C,D,E,F,O,O',G,G'}   \fill (\diem)circle(1.5pt);
        \end{tikzpicture}}
    \noindent $\dfrac{AG}{AO}=\dfrac{AG'}{AO'}=\dfrac{2}{3}$.\\
    Ta suy ra $GG' \parallel OO'$ và $OO' \parallel CE$ nên $GG' \parallel CE$.\\
    Mà $CE \subset \left(CDEF\right)$ nên $GG' \parallel \left(DCEF\right)$.
    }
\end{bt}

\begin{bt}%[1H4H2-2]
% nguon: Archive/SP_KhoLatex/ID6_old/1K4-12.tex #3286
%[Dự án 3 - TLDH - TeamTeXHoa - Võ Thị Thùy Trang]
    Cho tứ diện $ABCD$, gọi $G$ là trọng tâm tam giác $ABD$, $M$ là điểm trên cạnh $BC$ sao cho $MB=2MC$. Chứng minh $MG \parallel \left(ACD\right)$.
    \loigiai{
        \immini{
            Gọi $E$ là trung điểm của $AD$.\\
            Ta có $\dfrac{BG}{BE}=\dfrac{2}{3}$ (do $G$ là trọng tâm của tam giác $ABD$).\\
            Mà $\dfrac{BM}{BC}=\dfrac{2}{3}$ (do $MB=2MC$) \\\\
            Nên $\dfrac{BG}{BE}=\dfrac{BM}{BC}$. \\\\
            Suy ra $MG \parallel CE$.\\
            Mà $CE \subset \left( ACD \right)$ do đó $MG \parallel \left(ACD\right)$.}
        {   \begin{tikzpicture}[scale=0.7,>=stealth, font=\footnotesize, line join=round, line cap=round]
                \coordinate (D) at (2,4);
                \coordinate (A) at (0,0);
                \coordinate (B) at (1.5,-1.5);
                \coordinate (C) at (5,0);
                \coordinate[label=above:$E$] (E) at ($(A)!0.5!(D)$);
                \coordinate[label=below:$M$] (M) at ($(B)!2/3!(C)$);
                \coordinate[label=left:$G$] (G) at ($(B)!2/3!(E)$);
                \draw (A) node[left]{$A$}--(B) node[below]{$B$}--(C) node[right]{$C$}--(D) node[above]{$D$}--(A) (D)--(B)--(E);
                \draw[dashed] (A)--(C) (B)--(E)--(C) (G)--(M); 
                \foreach \diem in {A,B,C,D,E,M,G}   \fill (\diem)circle(1.5pt);
        \end{tikzpicture}}
    }
\end{bt}

\begin{bt}%[1H4H2-2]
% nguon: Archive/SP_KhoLatex/ID6_old/1K4-12.tex #3287
%[Dự án 3 - TLDH - TeamTeXHoa - Võ Thị Thùy Trang]
    Cho tứ diện $ABCD$. Gọi $M,N$ lần lượt là trọng tâm của các tam giác $ABC$ và $BCD$. Chứng minh rằng $MN \parallel \left(ABD\right), MN \parallel \left(ACD\right)$.
    \loigiai{
        \immini{
            Gọi $H$ là trung điểm của $BC$.\\
            Ta có $M \in AH, N \in DH$.\\
            Do đó $\dfrac{HM}{HA}=\dfrac{HN}{HD}=\dfrac{1}{3}$ (tính chất trọng tâm của tam giác)\\
            Nên $MN \parallel AD$. \\
            Khi đó\\
            $MN \parallel AD$ mà $AD \subset \left(ABD\right)$ nên $MN \parallel \left(ABD\right)$.\\
            $MN \parallel AD$ mà $AD \subset \left(ACD\right)$ nên $MN \parallel \left(ACD\right)$.}
        {   \begin{tikzpicture}[scale=0.7,>=stealth, font=\footnotesize, line join=round, line cap=round]
                \coordinate (D) at (2,4);
                \coordinate (A) at (0,0);
                \coordinate (B) at (1.5,-1.5);
                \coordinate (C) at (5,0);
                \coordinate[label=below:$H$] (H) at ($(C)!0.5!(B)$);
                \coordinate[label=below:$M$] (M) at ($(A)!2/3!(H)$);
                \coordinate[label=right:$N$] (N) at ($(D)!2/3!(H)$);
                \draw (A) node[left]{$A$}--(B) node[below]{$B$}--(C) node[right]{$C$}--(D) node[above]{$D$}--(A) (D)--(B) (D)--(H);
                \draw[dashed] (A)--(C) (A)--(H) (M)--(N); 
                \foreach \diem in {A,B,C,D,H,M,N}   \fill (\diem)circle(1.5pt);
        \end{tikzpicture}}  
    }
\end{bt}

\begin{bt}%[1H4H2-2]
% nguon: Archive/SP_KhoLatex/ID6_old/1K4-12.tex #3289
%[Dự án 3 - TLDH - TeamTeXHoa - Võ Thị Thùy Trang]
    Cho hình chóp $S.ABCD$ có đáy là $ABCD$ là hình bình hành; $F$, $G$ lần lượt là trung điểm của $AB$ và $CD$.
    \begin{enumerate}
        \item Chứng minh rằng $FG$ song song với các mặt phẳng $\left(SAD\right)$  và $\left( SBC\right)$ .
        \item Gọi $E$ là trung điểm của $SA$. Chứng minh rằng $SB, SC$ song song với mặt phẳng $\left(FGE\right)$ .
    \end{enumerate}
    \loigiai{
        \immini{
        \begin{enumerate}
            \item Ta có
            $\heva{&EF\parallel AD\\&AD\subset (SAD) }\Rightarrow EF \parallel (SAD)$.\\
            Tương tự ta có
            $\heva{&FG\parallel BC\\&BC\subset (SBC) }\Rightarrow FG \parallel (SBC)$.
            \item Gọi $H=SD\cap (EFG)$. Ta có
        \end{enumerate}}
    {   \begin{tikzpicture}[scale=0.35,>=stealth, font=\footnotesize, line join=round, line cap=round]  %hình chóp có đáy là hình thang.
            \coordinate[label=above left:$A$] (A) at (0,0);
            \coordinate[label=below left:$B$] (B) at (-4,-3);
            \coordinate[label=below left:$C$] (C) at (4,-3);
            \coordinate[label=right:$D$] (D) at (8,0);
            \coordinate[label=above:$S$] (S) at ($(A)+(-1,4)$);
            \coordinate[label=below:$F$] (F) at ($(B)!1/2!(A)$);
            \coordinate[label=below:$G$] (G) at ($(D)!1/2!(C)$);
            \coordinate[label=left:$E$] (E) at ($(A)!0.5!(S)$);
            \coordinate[label=above right:$H$] (H) at ($(S)!0.5!(D)$);
            \draw (S)--(B)--(C)--(D)--(S)--(C) (G)--(H);
            \draw[dashed] (A)--(B) (S)--(A)--(D) (H)--(E)--(F)--(G);
            \foreach \diem in {A,B,C,D,S,F,G,E,H}   \fill (\diem)circle(1.5pt);
    \end{tikzpicture}}
    $\heva{&(ABCD)\cap (EFG)=FG\\&(ABCD)\cap (SAD)=AD \\&(SAD)\cap (EFG)=EH\\& FG\parallel AD}\Rightarrow EH \parallel AD\parallel FG$.\\
Gọi $H$ là trung điểm của $SD$. Ta có   $\heva{&GH\parallel SC\\&GH\subset (EFG) }\Rightarrow SC \parallel (EFG)$.\\
Tương tự ta có $SB \parallel (EFG)$.
    }
\end{bt}

\begin{bt}%[1H4H3-5]
% nguon: Archive/SP_KhoLatex/ID6_old/1K4-12.tex #3290
%[Dự án 3 - TLDH - TeamTeXHoa - Võ Thị Thùy Trang]
    Cho hình chóp $S.ABCD$ có đáy $ABCD$ là hình bình hành.  $(\alpha)$ là mặt phẳng đi qua trung điểm $M$ của cạnh $SB$, song song với cạnh $AB$, cắt các cạnh $SA$, $SD$, $SC$ lần lượt tại $Q$, $P$ và $N$. Hãy xác định hình tính của tứ giác $MNPQ$?
    \loigiai{
        \immini{Ta có
            $\heva{&AB\parallel (\alpha)\\&M\in (\alpha)\cap (SAB)}\Rightarrow (\alpha)\cap (SAB)=MQ \parallel AB$.\\
            Mặt khác    $\heva{&DC\parallel AB\\&MQ\parallel AB}\Rightarrow DC \parallel QM \quad(1)$. \\
            Mà $QM\subset (\alpha)$ nên $DC\parallel (\alpha)$.\\
            Như vậy $\heva{&DC\parallel (\alpha)\\&PN= (\alpha)\cap (SCD)}\Rightarrow PN \parallel DC \quad(2)$.\\
            Từ (1) và (2) suy ra $MNPQ$ là hình thang.  }   
        {\begin{tikzpicture}[scale=0.35,>=stealth, font=\footnotesize, line join=round, line cap=round] %hình chóp có đáy là hình thang.
                \coordinate[label=below:$A$] (A) at (0,0);
                \coordinate[label=below left:$B$] (B) at (-4,-3);
                \coordinate[label=below left:$C$] (C) at (4,-3);
                \coordinate[label=right:$D$] (D) at (8,0);
                \coordinate[label=above:$S$] (S) at ($(A)+(-1,4)$);
                \coordinate[label=left:$M$] (M) at ($(B)!1/2!(S)$);
                \coordinate[label=right:$N$] (N) at ($(S)!1/2!(C)$);
                \coordinate[label=below:\tiny{Q}] (Q) at ($(A)!0.65!(S)$);
                \coordinate[label=above right:$P$] (P) at ($(D)!0.65!(S)$);
                \draw (S)--(B)--(C)--(D)--(S)--(C) (M)--(N)--(P);
                \draw[dashed] (A)--(B) (S)--(A)--(D) (M)--(Q)--(P);
                \foreach \diem in {A,B,C,D,S,M,N,P,Q}   \fill (\diem)circle(1.5pt);
        \end{tikzpicture}}
    }
\end{bt}

\begin{bt}%[1H4H3-2]
% nguon: Archive/SP_KhoLatex/ID6_old/1K4-12.tex #3291
%[Dự án 3 - TLDH - TeamTeXHoa - Võ Thị Thùy Trang]
    Cho hình chóp $S.ABCD$ có đáy là hình bình hành $ABCD$, tâm $O$. Gọi $M$ và $N$ lần lượt là trung điểm của $SA$ và $SD$.
    \begin{enumerate}
        \item Chứng minh $MN \parallel (SBC)$, $SB \parallel (OMN)$, $SC\parallel (OMN)$.
        \item Xác định thiết diện của hình chóp cắt bởi mặt phẳng $(OMN)$. Thiết diện là hình gì?
    \end{enumerate}
    \loigiai{\immini{
        \begin{enumerate}
            \item Ta có $MN\parallel AD$ ($MN$ là đường trung bình của tam giác $SAD$ ) và $AD\parallel BC$ (tứ giác $ABCD$ là hình bình hành), suy ra $MN\parallel BC$.\\
            Mà $BC\subset (SBC)$ nên $MN\parallel (SBC)$.\\
            Ta có $ON\parallel SB$ ($ON$ là đường trung bình của tam giác $(SBD)$ nên $ON \subset (OMN)$.\\
            Do đó $SB\parallel(OMN)$.\\
            Ta có $OM\parallel SC$ ($OM$ là đường trung bình của $\triangle{SAC})$ và $OM \subset(OMN)$.\\
            Vậy $SG\parallel (OMN)$.
            \item 
        \end{enumerate}
    }
{   \begin{tikzpicture}[scale=0.7,>=stealth, font=\footnotesize, line join=round, line cap=round]   %hình chóp có đáy là HCN.
        \coordinate[label=above right:$A$] (A) at (0,0);
        \coordinate[label=below left:$B$] (B) at (-2,-3);
        \coordinate[label=below left:$C$] (C) at (4,-3);
        \coordinate[label=right:$D$] (D) at (6,0);
        \coordinate[label=above:$S$] (S) at ($(A)+(0.5,4)$);
        \coordinate[label=below:$O$] (O) at ($(B)!.5!(D)$);
        \coordinate[label=right:$N$] (N) at ($(S)!0.5!(D)$);
        \coordinate[label=above right:$M$] (M) at ($(A)!0.5!(S)$);
        \coordinate[label=below right:$P$] (P) at ($(A)!0.5!(B)$);
        \coordinate[label=right:$Q$] (Q) at ($(C)!0.5!(D)$);
        \fill[color=gray!30] (O)--(N)--(M);
        \draw (S)--(B)--(C)--(D)--(S)--(C) (Q)--(N);
        \draw[dashed] (C)--(A)--(B) (S)--(A)--(D)--(B)  (Q)--(P)--(M)--(N) (M)--(O)--(N);
        \foreach \diem in {A,B,C,D,S,O,M,N,P,Q} \fill (\diem)circle(1.5pt);
    \end{tikzpicture}
}
Gọi $P$ và $Q$ lần lượt là trung điểm của $AB$ và $CD$. Từ đó có $PQ\parallel AD$, suy ra $PQ\parallel MN$.\\
Vậy $MN$ và $PQ$ đồng phẳng, nghĩa là $(OMN) \equiv(MNPQ)$.\\
Ta có thiết diện do $mp(OMN)$ cắt hình chóp là hình thang $MNPQ$ ($ MN\parallel PQ)$.
}
\end{bt}

\begin{bt}%[1H4H3-3]
% nguon: Archive/SP_KhoLatex/ID6_old/1K4-12.tex #3292
%[Dự án 3 - TLDH - TeamTeXHoa - Võ Thị Thùy Trang]
     Cho tứ diện $ABCD$. Gọi $I$ và $J$ lần lượt là trung điểm của $AB$ và $CD$, $M$ là một điểm trên đoạn
    $IJ$. Gọi $(P)$ là mặt phẳng qua $M$, song song với $AB$ và $CD$.
    \begin{enumerate}
        \item Tìm giao tuyến của mặt phẳng $(P)$ và mặt phẳng $(ICD)$.
        \item Xác định thiết diện của tứ diện với mặt phẳng $(P)$. Thiết diện là hình gì?
    \end{enumerate}
        \loigiai{\immini{
            \begin{enumerate}
                \item Ta có $\heva{&CD\parallel (P)\\& CD\subset(ICD) \\&M\in (P)\cap (ICD)}\\ \Rightarrow (P)\cap (ICD)=Mx\parallel CD$.\\
                Trong $mp(I C D)$ ta có $Mx$ cắt $IC$ tại $E$ và cắt $ID$ tại $F$. Suy ra $EF=(P) \cap(ICD)$.
                \item Ta có $\heva{&AB\parallel (P)\\& AB\subset(ABC) \\&E\in (P)\cap (ABC)}\\ \Rightarrow (P)\cap (ABC)=Ey\parallel AB$.\\
                Trong $mp(ABC)$ ta có $Ey$ cắt $BC$ tại $P$ và cắt $AC$ tại $S$. Suy ra $PS=(P) \cap(ABC)$.
            \end{enumerate}}
        {   \begin{tikzpicture}[scale=0.7, font=\footnotesize, line join=round, line cap=round, >=stealth]
                \coordinate (A) at (1.5,5);
                \coordinate (B) at (0,0);
                \coordinate (C) at (3.5,-2.5);
                \coordinate (D) at (7,0);
                \coordinate[label=left :$I$] (I) at ($(A)!0.55!(B)$);
                \coordinate[label=below :$J$] (J) at ($(D)!0.4!(C)$);
                \coordinate[label=above :$M$] (M) at ($(I)!0.55!(J)$);
                \coordinate[label=right :$E$] (E) at ($(I)!0.55!(C)$);
                \coordinate[label=below right :$F$] (F) at ($(I)!0.55!(D)$);
                \coordinate[label=below :$P$] (P) at ($(C)!0.45!(B)$);
                \coordinate[label=left:$S$] (S) at ($(C)!0.45!(A)$);    
                \coordinate[label=below :$Q$] (Q) at ($(D)!0.45!(B)$);
                \coordinate[label=above :$R$] (R) at ($(D)!0.45!(A)$);
                \draw (B) node[left]{$B$}--(C) node[below]{$C$}--(D) node[right]{$D$}--(A) node[above]{$A$}--(A)--(B) (A)--(C) (I)--(C) (P)--(S)--(R);
                %--(M) node[below right]{$M$}--(N) node[above right]{$N$} (S)--(B);
                \draw[dashed] (I)--(J) (B)--(D)--(I) (E)--(F) (R)--(Q)--(P);
                \foreach \diem in {A,B,C,D,I,J,M,P,Q,R,S,E,F}   \fill (\diem)circle(1.5pt);
        \end{tikzpicture}}
    Ta có $\heva{&AB\parallel (P)\\& AB\subset(ABD) \\&F\in (P)\cap (ABD)} \Rightarrow (P)\cap (ABD)=Ft\parallel AB$.\\
    Trong $mp(ABD)$ ta có $Ft$ cắt $BD$ tại $Q$ và cắt $AD$ tại $R$.
    Suy ra $QR=(P)\cap(ABD)$.\\
    Khi đó $PQ=(P)\cap(CBD)$ và $RS=(P) \cap(ACD)$.
    Vậy thiết diện cần tìm là tứ giác $PQRS$.\\
    Theo chứng minh trên ta có thể suy ra được $PS\parallel AB$,  $QR\parallel AB$ nên $PS\parallel QR\quad(1)$.\\
    Mặt khác, ta có $\heva{&\heva{&CD\parallel (P)\\&RS=(P)\cap (ACD)} \Rightarrow RS\parallel CD\\&\heva{&CD\parallel (P)\\&PQ=(P)\cap (BCD)} \Rightarrow PQ\parallel CD} \Rightarrow RS\parallel PQ\quad(2)$.\\
    Từ (1) và (2) suy ra thiết diện $PQRS$ là hình bình hành.       
    }
\end{bt}

\begin{bt}%[1H4H3-3]
% nguon: Archive/SP_KhoLatex/ID6_old/1K4-12.tex #3293
%[Dự án 3 - TLDH - TeamTeXHoa - Võ Thị Thùy Trang]
    Cho hình chóp $S.ABCD$. Gọi $M$, $N$ là hai điểm bất kì trên $SB$, $CD$. Mặt phẳng $(P)$ qua $MN$ và song song với $SC$.
    \begin{enumerate}
        \item Tìm các giao tuyến của $(P)$ với các mặt phẳng $(SBC)$, $(SCD)$, $(S A C)$.
        \item Xác định thiết diện của hình chóp với mặt phẳng $(P)$.
    \end{enumerate}
        \loigiai{
            \immini{
            \begin{enumerate}
                \item Trong mặt phẳng $(SBC)$, từ $M$ kẻ đường thẳng song song với $SC$ cắt $BC$ tại $Q$.\\
                Trong mặt phẳng $(SCD)$, từ $N$ kẻ đường thẳng song song với $SC$ cắt $SD$ tại $P$.\\
                Khi đó giao tuyến của $(P)$ với $(SBC)$ và $(SCD)$ lần lượt là $MQ$ và $NP$.\\
                Gọi $I=AC \cap NQ$. Từ $I$ kẻ đường thẳng song song với $SC$ cắt $SA$ tại $H$.\\
                Khi đó $(P) \cap(SAC)=IH$.
                \item Thiết diện của mặt phẳng $(P)$ với khối chóp là ngũ giác $MQNPH$.
            \end{enumerate} }
            {\begin{tikzpicture}[scale=0.7,>=stealth, font=\footnotesize, line join=round, line cap=round]  %hình chóp có đáy là HCN.
                    \coordinate[label=left:$A$] (A) at (-1,0);
                    \coordinate[label=below left:$B$] (B) at (2,-3);
                    \coordinate[label=below left:$C$] (C) at (5,-4);
                    \coordinate[label=right:$D$] (D) at (7,0);
                    \coordinate[label=above:$S$] (S) at ($(A)+(1.5,6)$);
                    \coordinate[label=left:$H$] (H) at ($(A)!.75!(S)$);
                    \coordinate[label=right:$N$] (N) at ($(C)!0.45!(D)$);
                    \coordinate[label=left:$M$] (M) at ($(B)!0.45!(S)$);
                    \coordinate[label=right:$P$] (P) at ($(S)!0.55!(D)$);
                    \coordinate[label=below:$Q$] (Q) at ($(B)!0.45!(C)$);
                    \path (intersection of C--A and Q--N) coordinate[label=above:$I$] (I);
                    \draw (S)--(B)--(C)--(D)--(S)--(C) (S)--(A)--(B) (M)--(Q) (N)--(P) (M)--(H);
                    \draw[dashed] (A)--(B) (S)--(A)--(D) (N)--(Q) (A)--(C) (I)--(H)--(P);
                    \foreach \diem in {A,B,C,D,S,M,N,Q,P,I,H}   \fill (\diem)circle(1.5pt);
            \end{tikzpicture}}
    }
\end{bt}

\begin{bt}%[1C1H5-4]
% nguon: Archive/SP_KhoLatex/ID6_old/1H1B5-4.tex #3352
%[HK1, Lê Quí Đôn - Đà Nẵng, 2018]%[Đức Nguyễn, 11EX1819]
Dựng về phía ngoài tam giác nhọn $ABC$ các tam giác đều $ABB'$, $ACC'$. Gọi $I$, $J$ lần lượt là trung điểm của $B'C$, $C'B$. Chứng minh rằng tam giác $AIJ$ đều.
	\loigiai{
	\begin{center}
		\begin{tikzpicture}[scale=1, font=\footnotesize,line join=round, line cap=round,>=stealth]
		\pgfmathsetmacro\a{4}
		\tkzDefPoint(0,0){B}
		\tkzDefShiftPoint[B](0:\a){C}
		\tkzDefShiftPoint[B](75:\a-0.5){A}
		\tkzDefCircle[radius](A,B)
		\tkzDefCircle[radius](B,A)
		\tkzInterCC(B,A)(A,B)    \tkzGetFirstPoint{B'}
		\tkzDefCircle[radius](A,C)
		\tkzDefCircle[radius](C,A)
		\tkzInterCC(A,C)(C,A)    \tkzGetFirstPoint{C'}
		\tkzDefMidPoint(B',C) \tkzGetPoint{I}
		\tkzDefMidPoint(B,C') \tkzGetPoint{J}
		\tkzDrawPolygon(B,B',A,C',C,B)
		\tkzDrawPolygon(A,I,J)
		\tkzDrawSegments(A,B A,C B',C B,C')
		\tkzLabelPoints[above](A,B',C',J)
		\tkzLabelPoints[below](B,C,I)
		\tkzDrawPoints[fill=black](A,B,C,B',C',I,J)
		\end{tikzpicture}
	\end{center}
Không mất tính tổng quát ta có thể giả sử tam giác $ABC$ có hướng như hình vẽ. Xét phép quay $\mathrm{Q}_{\left(A,\dfrac{\pi}{3}\right)}$ thì $B'\mapsto B$, $C\mapsto C'$, suy ra $B'C$ biến thành $BC'$ nên $I$ biến thành $J$.\\
Suy ra $AI=AJ$, $\widehat{IAJ}=\dfrac{\pi}{3}\Rightarrow $ tam giác $AIJ$ đều.
	}
\end{bt}

\begin{bt}%[1C1H6-4]
% nguon: Archive/SP_KhoLatex/ID6_old/1H1B5-4.tex #3353 
% KIEM_TRA: cần xem lại đáp số/dạng
%[HK1, Lê Quí Đôn - Đà Nẵng, 2018]%[Đức Nguyễn, 11EX1819]
	Cho nửa đường tròn đường kính $AB$. Gọi $C$ là điểm chuyển động trên nửa đường tròn và $I$ là điểm trên tia $AC$ sao cho $AI=BC$. Chứng minh $I$ di động trên một đường tròn cố định.
	\loigiai{
		\immini[0.05]{
			Không mất tính tổng quát, giả sử các điểm $A$, $B$, $C$ có hướng như hình vẽ trên. Gọi $M$ là điểm chính giữa của cung $\wideparen{AB}$, xét phép quay $\mathrm{Q}_{\left(M,-\dfrac{\pi}{2}\right)}$ thì $B\mapsto A$, đường thẳng $BC$ biến thành đường thẳng qua $A$ và vuông góc với $BC$, do $BC=AI$ nên $\mathrm{Q}_{\left(M,-\dfrac{\pi}{2}\right)}(C)=I$.\\
			Vì $C$ di động trên nửa đường tròn đường kính $AB$ nên $I$ chạy trên nửa đường tròn $(O')$ là ảnh của nửa đường tròn $(O)$ ($O$ là trung điểm của $AB$) qua $\mathrm{Q}_{\left(M,-\dfrac{\pi}{2}\right)}$.
		}{
			\begin{tikzpicture}[scale=1, font=\footnotesize,line join=round, line cap=round,>=stealth]
			\pgfmathsetmacro\a{4}
			\pgfmathsetmacro\r{\a/2}
			\tkzDefPoints{0/0/A}
			\tkzDefShiftPoint[A](0:\a){B}
			\tkzDefMidPoint(A,B) \tkzGetPoint{O}
			\tkzDefPointBy[rotation = center O angle 90](B) \tkzGetPoint{M}
			\tkzDrawArc(O,B)(A)
			\tkzDefPointBy[rotation = center O angle 45](B) \tkzGetPoint{C}
			\tkzCalcLength[cm](C,B) \tkzGetLength{rCB}
			\tkzCalcLength[cm](C,A) \tkzGetLength{rAC}
			\tkzDefPointBy[homothety = center A ratio \rCB/\rAC](C) \tkzGetPoint{I}
			\tkzDefCircle[circum](A,M,I)     \tkzGetPoint{O'}
			%\tkzDrawCircle[dashed](O',A)
			\tkzDefPointBy[rotation = center O' angle 180](A) \tkzGetPoint{A'}
			\tkzDrawArc(O',A)(A')
			\tkzDrawSegments(A,B B,C C,A A,M B,M M,I A,A')
			\tkzLabelPoints[left](O')
			\tkzLabelPoints[below](A,O,B,I)
			\tkzLabelPoints[above left](M)
			\tkzLabelPoints[right](C)
			\tkzDrawPoints[fill=black](A,B,C,O,M,O,O')
			\end{tikzpicture}
		}
	}
\end{bt}

\begin{bt}%[1H4H1-4]
% nguon: Archive/SP_KhoLatex/ID6_old/GK1.4.tex #3453

	Cho tứ diện $ABCD$ Gọi $K$ là trung điểm của cạnh $BC,{\rm \;}H$ là một điểm cố định trên cạnh $AC$ Mặt phẳng $(P)$ di động chứa $HK$, cắt các cạnh $BD$ và $AD$ lần lượt tại $M$ và $N$.	
	\begin{enumerate}
		\item  Giả sử cho trước điểm $M$ không là trung điểm của $BD$, hãy xác định điểm $N$
		\item  Tìm tập hợp giao điểm $I$ của hai đường $HM$ và $KN$ khi $M$ di động trên cạnh $BD$.
	\end{enumerate}
	\loigiai{
		\begin{center}
			\begin{tikzpicture}[scale=.7,line join=round, line cap=round,thick]
				\coordinate (A) at (0,4);
				\coordinate (B) at (-2,0);
				\coordinate (C) at (0,-3);
				\coordinate (D) at (3,0);
				\path ($(B)!0.5!(C)$) coordinate (K)
				($(A)!0.3!(C)$) coordinate (H)
				($(B)!0.7!(D)$) coordinate (M)
				(intersection of K--M and C--D) coordinate (E)
				(intersection of B--H and A--K) coordinate (F)
				(intersection of H--E and A--D) coordinate (N)
				(intersection of H--M and K--N) coordinate (I)
				;
				\draw(A)--(B)--(C)--(D)--cycle (A)--(C) (N)--(H)--(B) (A)--(K)--(H)--(D)--(E)--(H);
				\draw[dashed,thin](B)--(D) (M)--(N)--(K)--(E) (H)--(M) (K)--(D)--(F);
				\foreach \i/\g in {A/90,D/0,B/180,C/-90,K/-120,H/60,E/0,N/90,M/-90,I/-90,F/180}{\draw[fill](\i) circle (1.5pt) ($(\i)+(\g:3mm)$) node[scale=1]{$\i$};}
			\end{tikzpicture}
		\end{center}
		\begin{enumerate}
			\item  Trong (BCD): $KM\cap CD = \left\{E\right\}$.\\
			Trong (ACD): $HE\cap AD = \left\{N\right\}$.\\			
			Mà $HE\subset (P)$ nên $N=AD\cap (P)$ là điểm cần tìm.			
			\item  Ta có:
			$\left. \begin{array}{l}
				{I=HM\cap KN} \\
				{HM\subset (HBD)} \\
				{KN\subset (AKD)}
			\end{array}\right\}
			\Rightarrow I\in (HBD)\cap (AKD)\qquad (1)$\\			
			Trong (ABC): $BH\cap AK=\left\{F\right\}
			\Rightarrow F\in (HBD)\cap (AKD)$
			
			Mà $D\in (HBD)\cap (AKD)$ nên $DF=(HBD)\cap (AKD)$ (2)
			
			Từ (1) và (2) suy ra $I$ chạy trên đường thẳng cố định $DF$.
			
			\textit{Giới hạn:}
			
			Cho $M\to D$ thì $N\to D$ Khi đó $I\to D$
			
			Cho $M\to B$ thì $N\to A$ Khi đó $I\to F$
			
			Vậy tập hợp điểm $I$ là đoạn $DF$.
		\end{enumerate}
	}
	\end{bt}
